\chapter{预训练工程}
\label{lab:40_pre_training}

\noindent\textbf{配套文件：}\path{notebook/40_pre_training.ipynb}\par
\noindent\textbf{教材关联：}\bookxrefshort{ch:probability}、\bookxrefshort{ch:optimization-generalization}、\bookxrefshort{ch:pretraining-scaling}、\bookxrefshort{ch:training-stability}、\bookxrefshort{ch:distributed-training}。

\section*{实验目标}
验证词元目标、优化步和恢复状态的一致性。

\section*{准备及定位}
先阅读本册实验总则，确认Notebook中的依赖与资产单元。原理计算、库接口迁移和真实模型训练可能具有不同资源要求，应分别记录设备、版本及运行范围。

源码定位可从下列函数开始；应结合前置数据与配置单元阅读。
\begin{itemize}
\item \texttt{my\_pack\_documents}
\item \texttt{my\_collate}
\item \texttt{my\_build\_optimizer}
\end{itemize}

\section*{操作及观察}
\begin{enumerate}
\item 先检查文档打包边界、标签移位及有效词元计数。
\item 记录AdamW参数组和学习率计划，以相同数据顺序建立短程基线。
\item 保存检查点并恢复，比较接续训练与不中断训练的损失、参数及优化器状态。
\item 进入真实语料与Trainer路径前冻结数据和模型修订，记录训练预算及固定留出指标。
\end{enumerate}

\section*{结果判读及提交}
提交恢复对照、步数与词元预算、数据版本。仅重载权重不等于恢复训练；小规模真实训练也不能据此称为可用基座模型。

提交时附上所执行单元范围、环境与制品身份、原始结果及失败说明。将未运行项目明确列出，不能用Notebook中已有输出替代本次执行证据。
